%% =============================================================
%%  chapters/04_system_requirements.tex
%% =============================================================
\chapter{System Requirements and Design}
\label{ch:system-requirements}

\section{Use Case Definition and Target Scenarios}
\label{sec:use-case-definition}
The target system is a passenger car cabin monitoring function that
continuously estimates, for each seating position, whether the seat is
occupied, which passenger class occupies it (e.g. adult, or unknown), the occupant weight class, and the occupant's
body pose. Target scenarios include front and rear seating positions, a range of occupant
statures and clothing.


%% TODO: Exact target vehicle platform and number of monitored,
%% seating positions.

\section{Functional Requirements}
\label{sec:functional-requirements}
The system shall satisfy the following functional requirements:
\begin{enumerate}
    \item[\textbf{FR-01}] The system shall detect, per seating position,
    whether the seat is occupied or unoccupied.
    \item[\textbf{FR-02}] The system shall classify the seat occupant into
    predefined passenger classes (e.g. adult, or unknown).
    \item[\textbf{FR-03}] The system shall estimate the occupant's weight
    within a defined tolerance for adult.
    \item[\textbf{FR-04}] The system shall estimate the occupant's body
    pose for defined set  pose categories..
\end{enumerate}

%% TODO: Assign concrete numeric targets (e.g. exact class list, weight class) .

\section{Non-Functional Requirements}
\label{sec:non-functional-requirements}
The system shall additionally satisfy the following non-functional
requirements:
\begin{enumerate}
    \item[\textbf{NFR-01}] Occupancy detection accuracy shall exceed a
    defined target (e.g. 98\%) across the validation dataset described in
    Chapter~\ref{ch:data-acquisition}.
    \item[\textbf{NFR-02}] The system shall remain robust to sensor noise,
    partial occlusion, and seat cover variation without requiring
    per-vehicle recalibration.
    \item[\textbf{NFR-03}] All data collection, storage, and processing
    shall comply with the General Data Protection Regulation (GDPR) and
    applicable local data protection law.
    \item[\textbf{NFR-05}] The trained model shall degrade gracefully if
    one modality (camera or pressure) becomes temporarily unavailable.
\end{enumerate}

%% TODO: Define exact NFR-01 accuracy targets per output task and the
%% corresponding tolerance levels.

\section{System Architecture Overview}
\label{sec:system-architecture}
The proposed system architecture consists of four main stages: sensor
data acquisition (camera and pressure sensor array), preprocessing and
temporal synchronisation, the multi-task fusion model, and the output
interface that exposes per-seat occupancy, classification, weight, and
pose estimates to downstream vehicle systems. Figure
\ref{fig:system-architecture} illustrates this pipeline at block-diagram
level, showing the flow from raw sensor signals through synchronised
preprocessing into the shared fusion backbone and task-specific output
heads described in detail in Chapter~\ref{ch:fusion-architecture}.

\cite{mobihoc} This is a novel deep fusion architecture that can be used .

\begin{figure}[H]
\centering
%% TODO: Replace with the actual system architecture block diagram once
%% created (e.g. exported from draw.io or similar tool as PDF/PNG).
\includegraphics[width=0.85\textwidth]{system_architecture}
\caption{Block diagram of the proposed camera-pressure sensor fusion
system architecture, from raw sensor acquisition to multi-task output
estimation.}
\label{fig:system-architecture}
\end{figure}

\section{Design Decisions and Justification}
\label{sec:design-decisions}
Feature-level fusion was selected over data-level and decision-level
fusion because it allows each modality-specific branch to learn
representations suited to its own signal characteristics while still
enabling the network to learn cross-modal correlations before the final
task predictions are made, consistent with the fusion levels discussed in
Section~\ref{sec:fusion-levels}. A multi-task learning formulation was
chosen over four independent single-task models to exploit shared
structure between occupancy, classification, weight, and pose labels, and
to reduce the total inference compute footprint relative to running four
separate networks.


%% TODO: Add justification for specific backbone choice once architecture

